% The modulus: the operator bound on the fluctuation sector (S block).
% Lane: O (operator bridge), module YangMills/OS/SpatialSpectral.lean.
% REGIME: tricotomy; NO scores, NO desk language, NO commit chronology beyond
% the reproducibility freeze.
\documentclass[11pt]{article}
\usepackage[margin=1.1in]{geometry}
\usepackage{amsmath,amssymb,amsthm}
\usepackage{booktabs,longtable,array}
\usepackage[colorlinks=true,linkcolor=blue,urlcolor=blue,citecolor=blue]{hyperref}

\newtheorem{theorem}{Theorem}[section]
\newtheorem{lemma}[theorem]{Lemma}
\newtheorem{proposition}[theorem]{Proposition}
\newtheorem{corollary}[theorem]{Corollary}
\newtheorem{remark}[theorem]{Remark}
\theoremstyle{definition}
\newtheorem{definition}[theorem]{Definition}

\emergencystretch=3em
\tolerance=2000

\newcommand{\anchor}{d775f39f521cb18c9bad0cbba93fa97d6275bc56}
\newcommand{\osline}[3]{%
  \href{https://github.com/lluiseriksson/THE-ERIKSSON-PROGRAMME/blob/\anchor/YangMills/OS/#1.lean\#L#2}{\nolinkurl{#3}}}
\newcommand{\repofile}[2]{%
  \href{https://github.com/lluiseriksson/THE-ERIKSSON-PROGRAMME/blob/\anchor/#1}{\nolinkurl{#2}}}
\newcommand{\R}{\mathbb{R}}
\newcommand{\Z}{\mathbb{Z}}
\newcommand{\SU}{\mathrm{SU}}
\newcommand{\Om}{\Omega}

\title{The Modulus:\\
a Machine-Checked Operator Bound on the Fluctuation Sector\\
of the Coupled $\Z_2$ Slice\\
\large one construction, two debts --- and the one it does not pay}
\author{Lluis Eriksson}
\date{}

\begin{document}
\maketitle

\begin{abstract}
Two companion papers were left carrying the same debt from opposite sides.  The
gap paper proved that every eigenvalue of the coupled slice other than the
Perron eigenvalue is \emph{strictly} smaller in modulus, and said plainly that
this provides \emph{no modulus of separation} \cite{ErikssonGap}.  The bridge
paper proved that the Gibbs correlations of the spatial system are matrix
elements of a self-adjoint transfer operator, and obtained geometric decay only
\emph{under a contraction hypothesis it did not discharge} \cite{ErikssonGibbs}.
The missing step is identical in both: finitely many strict inequalities are not
an operator-norm bound.

\textbf{Proved.}  We construct $\mathrm{specGap}$, the largest $|\mu|$ over the
eigenvalues different from the Perron eigenvalue $\lambda$, and prove
$\mathrm{specGap} < \lambda$.  We then prove the operator bound: for every
observable $u$ orthogonal to the Perron vector,
$\|Ku\| \le \mathrm{specGap}\cdot\|u\|$.  That is exactly the hypothesis the
bridge paper carried, so its bound becomes unconditional.

\textbf{And a warning we state before anyone else has to.}  $\mathrm{specGap}
<\lambda$ is \emph{not} $\mathrm{specGap}<1$: the kernel is unnormalised, so
both are typically far above one and $\mathrm{specGap}^{N}$ \emph{grows}.  On
its own the unnormalised bound therefore controls growth, it does not exhibit
decay.  The rate that is below one is the \emph{relative} one,
$\mathrm{specRatio}=\mathrm{specGap}/\lambda<1$, and the decay statement is
that the fluctuation contribution is suppressed by $\mathrm{specRatio}^{N}$
\emph{relative to the Perron scale} $\lambda^{N}$.  Both forms are proved; only
the second is called decay.

The step that does not follow from the inequalities is the one about
eigenvectors \emph{at} $\lambda$: geometric simplicity, proved in the Perron
paper for an arbitrary eigenvector rather than a positive one
\cite{ErikssonPerron}, makes them \emph{invisible} to a fluctuation observable,
so the top term of the spectral sum vanishes instead of merely being bounded.

\textbf{Not proved.}  $\mathrm{specGap}$ \emph{depends on the extent}, and
nothing here bounds it away from $\lambda$ uniformly.  Direct diagonalisation
gives $\mathrm{specGap}/\lambda = 0.9205, 0.9829, 0.9964, 0.9992$ at
$L = 2,3,4,5$ for one parameter pair: a geometric bound whose rate tends to $1$
is empty in the limit, and that is reported, not hidden.

\textbf{And a correction to our own previous scope claim.}  An earlier version
said that bounding the partition function below requires identifying the index
of the top eigenvalue.  That was \emph{wrong}, as an external reading pointed
out.  Splitting the dressed constant observable along the Perron direction makes
both cross terms vanish by symmetry alone, and the resulting lower bound turns
the relative statement into a bound on the \emph{normalised} Gibbs two-point
function.  Both are now proved.  Reflection
positivity is untouched, and nothing in this paper is a claim about $\SU(N)$,
the continuum limit, or the Yang--Mills mass gap.
\end{abstract}

\section{The same debt, from two sides}\label{sec:debt}

The gap paper's headline is a strict separation without a modulus.  The bridge
paper's decay theorem carries a contraction hypothesis.  It is worth stating
precisely why these are one debt and not two.

A strict inequality $|\mu| < \lambda$ for each of finitely many eigenvalues does
give $\max|\mu| < \lambda$ --- taking a maximum over a finite set is not the
difficulty.  The difficulty is that $\max|\mu|$ over the non-Perron eigenvalues
is not, by itself, known to \emph{dominate the operator} on the orthogonal
complement of the Perron vector.  Getting from one to the other requires the
spectral decomposition: an orthonormal basis of eigenvectors, so that an
observable can be expanded and its image measured coordinate by coordinate.

That is what this paper supplies, and it pays both debts with one construction.
We prove the inequality $\|Ku\|\le\mathrm{specGap}\|u\|$; we do \emph{not}
prove that $\mathrm{specGap}$ \emph{equals} the restricted operator norm.  The
numerical section reports that the bound is attained, but attainment is
\textbf{verified}, not proved, and no theorem here depends on it.

\section{The construction}\label{sec:constr}

Let $K$ be a strictly positive symmetric kernel on a finite nonempty type, with
Perron data $(v,\lambda)$: $v > 0$ pointwise and $Kv = \lambda v$.

\begin{definition}\label{def:gap}
Let $\mu_j$ be the eigenvalues of $K$ and set
\[
  \mathrm{specGap} \;=\; \max_j
  \begin{cases}
    0 & \text{if } \mu_j = \lambda,\\
    |\mu_j| & \text{otherwise.}
  \end{cases}
\]
\end{definition}

\osline{SpatialSpectral}{172}{specGap}.
Eigenvalues equal to $\lambda$ contribute $0$ rather than $\lambda$.  That is
legitimate --- and it is the whole content of §\ref{sec:invisible}.

\begin{theorem}[the modulus]\label{thm:gaplt}
$\mathrm{specGap} < \lambda$.
\end{theorem}

\osline{SpatialSpectral}{184}{specGap\_lt}.
Each term of the maximum is either $0 < \lambda$, or $|\mu_j| < \lambda$ by the
gap paper applied to the $j$-th basis vector, which is a nonzero eigenvector
because it has unit norm.

\section{Why eigenvectors at the top may be discarded}\label{sec:invisible}

\begin{theorem}\label{thm:invisible}
If $\mu_j = \lambda$ and $u$ is orthogonal to the Perron vector, then $u$ has no
component along the $j$-th eigenvector.
\end{theorem}

\osline{SpatialSpectral}{205}{inner\_specBasis\_eq\_zero\_of\_top}.

\begin{remark}
This is the step that does not come from the inequalities, and it is why
Definition~\ref{def:gap} may assign $0$ to the top of the spectrum.  The Perron
paper proves geometric simplicity for an \emph{arbitrary} eigenvector with
eigenvalue $\lambda$, not merely a positive one \cite{ErikssonPerron}; hence the
$j$-th basis vector is a multiple of $v$, and orthogonality to $v$ kills the
coordinate outright.  Had only the positive case been available, the top term
could have been bounded but not removed, and the maximum in
Definition~\ref{def:gap} would have had to include $\lambda$ --- making
Theorem~\ref{thm:gaplt} false and the whole construction vacuous.
\end{remark}

\section{The operator bound}\label{sec:bound}

\begin{theorem}[the bound]\label{thm:bound}
For every $u$ with $\langle v,u\rangle = 0$,
\[
  \|Ku\| \;\le\; \mathrm{specGap}\cdot\|u\| .
\]
\end{theorem}

\osline{SpatialSpectral}{263}{norm\_act\_le\_specGap}.

The proof is the expansion in the eigenbasis.  Writing $c_j$ for the coordinates
of $u$, the coordinates of $Ku$ are $\mu_jc_j$
(\osline{SpatialSpectral}{223}{inner\_specBasis\_act}), so by Parseval
(\osline{SpatialSpectral}{245}{eucNorm\_sq\_eq\_sum})
\[
  \|Ku\|^2=\sum_j \mu_j^2c_j^2,\qquad \|u\|^2=\sum_j c_j^2 ,
\]
and each term is dominated: if $\mu_j=\lambda$ then $c_j=0$ by
Theorem~\ref{thm:invisible}, and otherwise $|\mu_j|\le\mathrm{specGap}$ by
definition.

\section{The endpoint}\label{sec:endpoint}

The bridge paper's fluctuation sector is invariant under a kernel that
\emph{fixes} the vacuum.  Here the kernel is unnormalised and the vacuum is an
eigenvector with eigenvalue $\lambda$, so invariance is reproved in that form
(\osline{SpatialSpectral}{356}{perp\_invariant\_eigen}); the scalar drops out
because the pairing is bilinear.  Iterating gives
\osline{SpatialSpectral}{371}{iterate\_norm\_le\_specGap}.

\begin{theorem}[the bridge bound, without its hypothesis]\label{thm:endpoint}
For the coupled spatial system, at every finite extent $L$, every $\beta$, and
every strictly positive source weight $w$: if the dressing of $A$ is orthogonal
to the Perron vector, then for every $N$
\[
  \Bigl|\sum_X A(X_0)A(X_N)\,\mathcal{W}(X)\Bigr|
  \;\le\; \mathrm{specGap}^{\,N}\;\bigl\|\sqrt{w}A\bigr\|^{2},
  \qquad \mathrm{specGap}<\lambda .
\]
\end{theorem}

\osline{SpatialSpectral}{404}{gibbs\_pathSum\_bound\_unconditional}.

\begin{remark}[why that is not yet decay]\label{rem:notdecay}
$\mathrm{specGap}<\lambda$ does not give $\mathrm{specGap}<1$.  The kernel is
unnormalised: in the probe of §\ref{sec:verified}, $L=2$ with $\beta=0.8$,
$\gamma=1.2$ gives $\lambda\approx56.89$ and
$\mathrm{specGap}\approx52.37$, so $\mathrm{specGap}^{N}$ grows.
Theorem~\ref{thm:endpoint} is a \emph{sub-Perron growth} bound.  The decay is
Theorem~\ref{thm:reldecay}.
\end{remark}

\begin{theorem}[the decay statement]\label{thm:reldecay}
With $\mathrm{specRatio}=\mathrm{specGap}/\lambda$, under the same hypotheses,
\[
  \Bigl|\sum_X A(X_0)A(X_N)\,\mathcal{W}(X)\Bigr|
  \;\le\; \mathrm{specRatio}^{\,N}\;\Bigl(\lambda^{N}
  \bigl\|\sqrt{w}A\bigr\|^{2}\Bigr),
  \qquad \mathrm{specRatio}<1 .
\]
Equivalently, the normalised kernel $\lambda^{-1}K$ contracts the fluctuation
sector by $\mathrm{specRatio}<1$.
\end{theorem}

\osline{SpatialSpectral}{445}{gibbs\_pathSum\_relative\_decay},
\osline{SpatialSpectral}{312}{specRatio\_lt\_one},
\osline{SpatialSpectral}{335}{norm\_act\_normalizedKernel\_le}.

\begin{remark}[non-vacuity]
Dressing is multiplication by a strictly positive function, so every fluctuation
vector is the dressing of an observable: whenever the configuration space has
two points there is a \emph{nonzero} $A$ satisfying the hypothesis, and
Theorem~\ref{thm:endpoint} is not quantified over an empty set
(\osline{SpatialSpectral}{471}{exists\_nonzero\_dressed\_fluctuation}).
\end{remark}

\section{Scope: three things this does not give}\label{sec:scope}

\paragraph{Nothing uniform in the extent.}  This is the substantive limitation.
$\mathrm{specGap}$ is built from the spectrum at a fixed extent and depends on
it.  Direct dense diagonalisation gives, for $\beta=0.8$, $\gamma=1.2$,
\[
  \mathrm{specGap}/\lambda \;=\; 0.6640,\;0.9205,\;0.9829,\;0.9964,\;0.9992
  \qquad\text{at } L=1,\dots,5,
\]
and for $\beta=0.3$, $\gamma=0.4$,
\[
  0.2913,\;0.4665,\;0.5518,\;0.5938,\;0.6159 .
\]
The first family collapses towards $1$, where a bound of the form
$\mathrm{specGap}^N$ says nothing in the limit.  These numbers are
\textbf{verified} and \textbf{not proved}, no theorem depends on them, and they
are printed because a paper that reported only Theorem~\ref{thm:endpoint} would
be reporting the half that does not settle the question.

\paragraph{A bound, not an identification.}  We prove
$\|Ku\|\le\mathrm{specGap}\|u\|$ on the fluctuation sector.  We do \emph{not}
prove that $\mathrm{specGap}$ equals the restricted operator norm.  The probe
reports attainment as \textbf{verified}, and §\ref{sec:denominator} explains
why we now expect this to be cheap rather than blocked.

\paragraph{Still not clustering in the infinite-volume sense.}
Theorem~\ref{thm:corr} bounds the normalised two-point function at each fixed
extent.  It says nothing about the limit: the rate is
$\mathrm{specRatio}(L)$, and the measured evidence is that it tends to $1$.
The word \emph{clustering} is reserved for a statement that survives that limit,
and no such statement is made.

\paragraph{No claim about the physical target.}  Reflection positivity is
untouched, and nothing here concerns $\SU(N)$, the continuum limit, or the Clay
problem \cite{Clay}.  The distance to that problem is unchanged.

\section{The denominator, and a retraction}\label{sec:denominator}

The first version of this paper asserted, in its scope section, that a lower
bound on the partition function needs the index of the top eigenvalue
identified.  \textbf{That was wrong.}  An external reading supplied the argument
below, and it uses no index at all.

Let $b$ be the dressed constant observable, $\Om$ the unit Perron vacuum,
$c=\langle\Om,b\rangle$ and $u=b-c\,\Om$, so that $u\perp\Om$.  Then
\[
  \langle b,K^{N}b\rangle
  \;=\; c^{2}\lambda^{N}+\langle u,K^{N}u\rangle ,
\]
because \emph{both} cross terms vanish: $\langle\Om,K^{N}u\rangle
=\langle K^{N}\Om,u\rangle=\lambda^{N}\langle\Om,u\rangle=0$ by symmetry and
the eigen-equation, and the other by the same computation.  Nothing here asks
which basis vector carries $\lambda$.

\osline{SpatialSpectral}{558}{quadForm\_split}.

\begin{theorem}[the partition function from below]\label{thm:denom}
$\langle b,K^{N}b\rangle \;\ge\; c^{2}\lambda^{N}
 -\mathrm{specGap}^{\,N}\|u\|^{2}$, and the left-hand side is the partition
function.
\end{theorem}

\osline{SpatialSpectral}{596}{quadForm\_lower}.
Since $c>0$ (both $\Om$ and $b$ are strictly positive) and
$\mathrm{specRatio}<1$, the right-hand side is positive for all large $N$.

\begin{theorem}[the normalised two-point function]\label{thm:corr}
For any $B$ bounding the unnormalised sum and any $D>0$ bounding the partition
function below,
\[
  \bigl|\mathbb{E}[A(X_0)A(X_N)]\bigr|\;\le\;B/D .
\]
Combining Theorem~\ref{thm:endpoint} with Theorem~\ref{thm:denom} gives a bound
of the form $\mathrm{const}\cdot\mathrm{specRatio}^{\,N}$ at fixed extent.
\end{theorem}

\osline{SpatialSpectral}{652}{gibbsCorr\_bound\_of\_partition\_lower}.

\begin{remark}
The same reading observes that the \emph{sharpness} of $\mathrm{specGap}$ is
likewise not blocked by the missing index: the maximum defining it is attained
at some basis vector, that vector is orthogonal to the Perron vector because
eigenvectors with different eigenvalues are, and it realises the bound.  We have
not formalised this, and say so rather than claim it.
\end{remark}

\section{What is verified and not proved}\label{sec:verified}

A machine-checked proof guarantees that the definitions are internally
consistent, not that they name the object the prose claims.  The construction was
therefore recomputed by direct dense diagonalisation, by a script with no access
to the formalisation.  Three things were checked: that
$\mathrm{specGap} < \lambda$ in every case; that
$\|Ku\|\le \mathrm{specGap}\|u\|$ on $400$ random fluctuation vectors per case;
and that the bound is \emph{attained} by the subdominant eigenvector, so it is
sharp rather than a loose over-estimate.  All cases pass.

One consequence is worth recording.  The measured $\mathrm{specGap}/\lambda$
reproduces, digit for digit, the subdominant ratios the gap paper reported as
measured and unnamed \cite{ErikssonGap}.  The object constructed here is that
number.

\section{Formalisation notes}\label{sec:formal}

\paragraph{The library had the hard half.}  mathlib provides an orthonormal
eigenbasis and real eigenvalues for a Hermitian matrix in finite dimension.  The
work in this module is not the spectral theorem; it is the bridging --- writing
the kernel as a Hermitian matrix, matching mathlib's eigen-equation to the form
the companion papers state, and translating norms and pairings between the two
spellings.  That bridging is where the effort went, and it is the usual place.

\paragraph{Where the unnormalised kernel forced a reproof.}  The bridge module's
invariance lemma assumes the vacuum is \emph{fixed}.  Its kernel is
unnormalised, so the vacuum is an eigenvector with eigenvalue $\lambda\neq1$ in
general, and the lemma does not apply verbatim.  Rather than reopen a frozen
module, the invariance is reproved here in the eigenvector form.  It is three
lines, because the eigenvalue factors out of a bilinear pairing.

\section{Reproducibility}\label{sec:repro}

All artifacts are at source checkpoint \texttt{\anchor} on branch \texttt{main}.
Every hyperlink was resolved against that commit, and checked to point at an
actual declaration line, before compilation.

\begin{center}
\begin{tabular}{ll}
\toprule
Quantity & Value \\
\midrule
Full core build & 8430 jobs, success \\
Oracle commands & 2665 \\
\quad with axiom dependencies & 2642 \\
\quad axiom-free & 23 \\
Distinct axioms across all outputs & \texttt{propext}, \texttt{Quot.sound}, \texttt{Classical.choice} \\
Occurrences of \texttt{sorryAx} & 0 \\
Project axioms & 0 \\
Errors & 0 \\
\bottomrule
\end{tabular}
\end{center}

The module is
\repofile{YangMills/OS/SpatialSpectral.lean}{YangMills/OS/SpatialSpectral.lean};
the probe is
\repofile{scripts/probe\_spatial\_spectral.py}{scripts/probe\_spatial\_spectral.py};
verdicts are in
\repofile{docs/VERIFICATION-LEDGER.md}{docs/VERIFICATION-LEDGER.md}.

\begin{thebibliography}{9}

\bibitem{ErikssonPerron}
L.~Eriksson,
\emph{The Vacuum Was Never Absent: A Machine-Checked Perron Theorem for
Strictly Positive Kernels, and the Coupled-Slice Vacuum at Every Spatial
Extent}, viXra, 2026.

\bibitem{ErikssonGap}
L.~Eriksson,
\emph{Strict but Not Uniform: a Machine-Checked Spectral Gap at Every Finite
Extent of the Coupled Slice}, viXra, 2026.

\bibitem{ErikssonGibbs}
L.~Eriksson,
\emph{The Measure the Spectral Results Were About: a Machine-Checked Transfer
Bridge for the Spatial $\Z_2$ Slice}, viXra, 2026.

\bibitem{HornJohnson}
R.~A.~Horn and C.~R.~Johnson,
\emph{Matrix Analysis}, 2nd ed., Cambridge Univ.\ Press, 2013.
(Spectral theorem for Hermitian matrices; the operator norm of a self-adjoint
matrix as the spectral radius.)

\bibitem{Seneta}
E.~Seneta,
\emph{Non-negative Matrices and Markov Chains}, 2nd ed., Springer, 1981.
(Perron--Frobenius theory, simplicity of the dominant eigenvalue, and the
subdominant ratio as the rate of convergence.)

\bibitem{Baladi}
V.~Baladi,
\emph{Positive Transfer Operators and Decay of Correlations},
World Scientific, 2000.
(Transfer operators, the spectral gap, and decay of correlations relative to the
leading eigenvalue.)

\bibitem{Simon}
B.~Simon,
\emph{The Statistical Mechanics of Lattice Gases, Volume I},
Princeton Univ.\ Press, 1993.
(Transfer-matrix formulation of lattice Gibbs measures and correlation decay.)

\bibitem{Lean4}
L.~de~Moura and S.~Ullrich,
\emph{The Lean 4 Theorem Prover and Programming Language},
CADE 28 (2021), 625--635.

\bibitem{Mathlib}
The mathlib Community,
\emph{The Lean Mathematical Library},
CPP 2020, 367--381.

\bibitem{Clay}
A.~Jaffe and E.~Witten,
\emph{Quantum Yang--Mills Theory},
Clay Mathematics Institute Millennium Problem Description, 2000.

\end{thebibliography}

\end{document}
